\section{Motivation}
\label{sec:motivation}

Requirements Engineering (RE) plays a foundational role in software engineering by transforming high-level stakeholder intent into precise technical specifications~\cite{sommerville2015software}. However, initial stakeholder intent is rarely complete; it is often vague, fragmented, and filled with unstated assumptions. When these ambiguities are left unaddressed during early RE stages, they propagate into downstream implementation and testing, resulting in system-level defects and expensive late-stage rework. Recently, the emergence of Large Language Models (LLMs) has enabled intent-driven software development, offering new opportunities to automate requirement elicitation and formalization.

To automate RE tasks, recent works have introduced LLM-based multi-agent frameworks that structure engineering activities into multi-step collaborative workflows~\cite{mare2024, jin2025iredev}. By partitioning responsibilities among specialized roles, such as analysts, modelers, and checkers, these approaches effectively automate requirement elicitation, artifact modeling, and quality negotiation~\cite{mare2024, cheng2026quare}. However, existing systems primarily focus on procedural workflow execution and system-centric role division. They often rely on generic or functional agent prompts rather than modeling the rich diversity of actual human stakeholders. Consequently, current multi-agent RE frameworks still struggle to fully surface implicit user expectations, domain-specific nuances, and edge cases that naturally emerge from multi-perspective stakeholder interactions.

Personas are fictional yet realistic archetypes of target end-users, first introduced by Alan Cooper~\cite{cooper2004inmates}. By consolidating users who share similar goals, motivations, and behaviours into concrete characters, personas keep actual users at the centre of the RE process and help reduce assumption bias, i.e., defining requirements based on what users want rather than on what engineers assume they want~\cite{karolita2023personas}. Accordingly, personas have been applied predominantly during requirements elicitation and analysis, where they support understanding end users, identifying requirements, and communicating with stakeholders~\cite{karolita2023personas}. Compared with generic role-based models, refining abstract roles into diverse, concrete personas exposes a broader spectrum of user needs and behaviours early in the requirements process~\cite{sim2014empowering}.

In practice, however, constructing personas remains costly and unreliable: it requires data collection from representative users, suffers from the lack of common standards, and adds time and effort to the RE process~\cite{karolita2023personas}; adoption is consequently uneven across organizations, and human aspects are often left out of persona descriptions~\cite{wang2025personas}. Large Language Models offer a promising alternative, as LLM-based agents can role-play diverse user personas to elicit requirements at scale~\cite{ataei2025elicitron}. Yet recent evidence also cautions that LLM-generated personas can introduce systematic biases in downstream tasks~\cite{li2026persona}, and that the persona fidelity of LLMs varies considerably across models~\cite{samuel2025personagym}. These observations motivate grounding multi-agent requirements elicitation in carefully constructed, domain-grounded personas.
